% !TEX root = ../main.tex
% =====================================================================
% Chapter 5 — Golden Bridge: Fibonacci-Lucas Generating Functions
% Lane: L5 (THEORY, R7 N/A, R14 N/A per trios#265 §2.2)
% Agent: perplexity-computer-l5-bridge
% Mission: trios#265 ONE SHOT — Flos Aureus PhD monograph
% R3 (>=1500 lines, >=2 cites, >=1 theorem with proof+qed) · R12 Lee/GVSU
% =====================================================================

\chapter{The Golden Bridge: Fibonacci-Lucas Generating Functions}


% Chapter Anchor header (Phase 1 UNIFY task 1.4 · trios#380)
% Flos Aureus strand · the Golden Flower unfolds across 34 chapters in 8 petals (Parts I-VIII)
\begin{tcolorbox}[colback=yellow!4,colframe=yellow!50!brown,title={\textbf{Flos Aureus} \textbf{FA.05 Golden Bridge}}]
  \textbf{Petal:} Part II --- \emph{The Expansion} \\
  \textbf{Anchor:} \(\varphi^{2} + \varphi^{-2} = 3\) (Trinity Identity, INV-22) \\
  \textbf{Motif:} Fibonacci–Lucas generating functions \\
  \textbf{Lane:} L5 (Flos Aureus strand) \\
  \textbf{Theorems in chapter:} 46 \\
  \textbf{Coq link:} \filepath{trinity-clara/proofs/igla/lucas\_closure\_gf16.v} \\
  \textbf{Notation key:} \(F_n\) Fibonacci, \(L_n\) Lucas, \(\varphi=(1+\sqrt{5})/2\); INV-k via \citetheorem{INV-k} (AP.F)
\end{tcolorbox}

\begin{figure}[H]
\centering
\makebox[\linewidth][c]{\includegraphics[width=1.18\linewidth,keepaspectratio]{11-vesica-piscis.jpg}}
\end{figure}


\label{fa_05:ch:golden-bridge}

\epigraph{
The single rational function $1/(1 - x - x^2)$ is the generating
function of the Fibonacci sequence; the same denominator carries the
Lucas sequence; and the analytic-arithmetic bridge between them is the
content of this chapter.
}{Lee/GVSU style preamble.}

\section{Introduction and Statement of the Main Result}
\label{fa_05:sec:05-intro}

The Trinity Anchor identity $L_2 = 3$ has, by Chapter~\ref{ch:lucas-ring}
(L6), an algebraic substrate in the Lucas ring $\mathcal{L}=\mathbb{Z}[\varphi]$,
and by Chapter~\ref{ch:lucas-ladder} (L4), a recurrence-theoretic
substrate in the Lucas ladder $\{L_n\}$. The present chapter establishes
the analytic bridge between these two substrates: the generating
functions
\[
  F(x) = \sum_{n \ge 0} F_n x^n = \frac{x}{1 - x - x^2}
  \quad\text{and}\quad
  L(x) = \sum_{n \ge 0} L_n x^n = \frac{2 - x}{1 - x - x^2}.
\]

Both generating functions are rational, share the denominator
$1 - x - x^2$, and admit partial-fraction decompositions over the
quadratic extension $\mathbb{Q}(\varphi)$ that recover the Binet formula
proven in Chapter~\ref{ch:lucas-ladder}. The aim of this chapter is to
prove the following structural theorem.

\begin{theorem}[Golden-Bridge Structure Theorem]
\label{thm:05-bridge}
The Fibonacci and Lucas generating functions $F(x), L(x) \in \mathbb{Q}(x)$
satisfy the following identities simultaneously:
\begin{enumerate}
  \item $F(x) = x / (1 - x - x^2)$ and $L(x) = (2 - x)/(1 - x - x^2)$
        (Theorem~\ref{thm:05-genfn-closed}).
  \item Both functions admit the partial-fraction decomposition over
        $\mathbb{Q}(\varphi)$:
        $F(x) = \frac{1}{\sqrt{5}}\left(\frac{1}{1 - \varphi x} - \frac{1}{1 - \psi x}\right)$,
        $L(x) = \frac{1}{1 - \varphi x} + \frac{1}{1 - \psi x}$
        (Theorem~\ref{thm:05-partial-frac}).
  \item Their radius of convergence is $|x| < 1/\varphi$, and on this
        disc both $F(x)$ and $L(x)$ are analytic
        (Theorem~\ref{thm:05-radius}).
  \item $L(x) - 2 = x \cdot F(x) + x^2 \cdot L(x)$ (Lucas-Fibonacci coupling
        identity, Theorem~\ref{thm:05-coupling}).
  \item The coefficient of $x^2$ in $L(x)$ is $L_2 = 3$, the Trinity Anchor
        (Theorem~\ref{thm:05-anchor-as-coeff}).
\end{enumerate}
\end{theorem}

\begin{proof}[Outline]
The five clauses are proved in §\ref{sec:05-closed-form}–§\ref{sec:05-anchor-coeff}.
Strand I (§\ref{sec:05-strand-i}) is the algebraic-rational generating-function
content. Strand II (§\ref{sec:05-strand-ii}) is the analytic content
(radius of convergence, asymptotics, Hankel). Strand III (§\ref{sec:05-strand-iii})
is the combinatorial-arithmetic bridge to L4 and L6. \qed
\end{proof}

The chapter is organised as follows. Section~\ref{sec:05-prelim}
fixes notation and recalls the recurrences for $F_n, L_n$. Sections
\ref{sec:05-closed-form}, \ref{sec:05-partial-frac}, \ref{sec:05-radius},
\ref{sec:05-coupling}, and \ref{sec:05-anchor-coeff} prove the five
clauses of Theorem~\ref{thm:05-bridge}. Sections~\ref{sec:05-strand-i},
\ref{sec:05-strand-ii}, \ref{sec:05-strand-iii} develop the three
strands of consequences. Section~\ref{sec:05-falsification} discusses
falsification. Appendices A through Z record auxiliary computations.

\section{Preliminaries: Recurrences and Initial Conditions}
\label{sec:05-prelim}

\begin{definition}
\label{fa_05:def:05-fib}
The Fibonacci sequence $\{F_n\}_{n \ge 0}$ is defined by
$F_0 = 0$, $F_1 = 1$, $F_{n+2} = F_{n+1} + F_n$ for $n \ge 0$.
\end{definition}

\begin{definition}
\label{fa_05:def:05-luc}
The Lucas sequence $\{L_n\}_{n \ge 0}$ is defined by
$L_0 = 2$, $L_1 = 1$, $L_{n+2} = L_{n+1} + L_n$ for $n \ge 0$.
\end{definition}

\begin{lemma}
\label{fa_05:lem:05-fib-vals}
The first ten Fibonacci numbers are
$F_0 = 0, F_1 = 1, F_2 = 1, F_3 = 2, F_4 = 3, F_5 = 5, F_6 = 8,
F_7 = 13, F_8 = 21, F_9 = 34$.
\end{lemma}

\begin{proof}
By repeated application of the recurrence $F_{n+2} = F_{n+1} + F_n$. \qed
\end{proof}

\begin{lemma}
\label{fa_05:lem:05-luc-vals}
The first ten Lucas numbers are
$L_0 = 2, L_1 = 1, L_2 = 3, L_3 = 4, L_4 = 7, L_5 = 11, L_6 = 18,
L_7 = 29, L_8 = 47, L_9 = 76$.
\end{lemma}

\begin{proof}
By repeated application of the recurrence $L_{n+2} = L_{n+1} + L_n$. \qed
\end{proof}

\begin{remark}
The number $L_2 = 3$ is the Trinity Anchor, the central object of the
Flos Aureus monograph. Throughout this chapter we trace the identity
$L_2 = 3$ as it appears in the generating function machinery.
\end{remark}

\section{Closed-Form Generating Functions (Clause 1)}
\label{sec:05-closed-form}

We now prove the closed-form expressions for $F(x)$ and $L(x)$.

\begin{theorem}
\label{thm:05-genfn-closed}
The formal power series
\[
  F(x) = \sum_{n \ge 0} F_n x^n
  \quad\text{and}\quad
  L(x) = \sum_{n \ge 0} L_n x^n
\]
admit the closed-form rational expressions
\[
  F(x) = \frac{x}{1 - x - x^2}
  \quad\text{and}\quad
  L(x) = \frac{2 - x}{1 - x - x^2}.
\]
\end{theorem}

\begin{proof}
\emph{Fibonacci case.} Multiply the recurrence $F_{n+2} = F_{n+1} + F_n$
by $x^{n+2}$ and sum over $n \ge 0$:
\[
  \sum_{n \ge 0} F_{n+2} x^{n+2} = x \sum_{n \ge 0} F_{n+1} x^{n+1} + x^2 \sum_{n \ge 0} F_n x^n.
\]
The left-hand side is $F(x) - F_0 - F_1 x = F(x) - x$. The right-hand
side is $x (F(x) - F_0) + x^2 F(x) = x F(x) + x^2 F(x)$. Hence
$F(x) - x = x F(x) + x^2 F(x)$, i.e., $F(x)(1 - x - x^2) = x$, and
$F(x) = x / (1 - x - x^2)$.

\emph{Lucas case.} Similarly, multiply $L_{n+2} = L_{n+1} + L_n$ by
$x^{n+2}$ and sum:
$L(x) - L_0 - L_1 x = x (L(x) - L_0) + x^2 L(x)$.
$L(x) - 2 - x = x L(x) - 2x + x^2 L(x)$.
$L(x) - 2 - x + 2x = x L(x) + x^2 L(x)$.
$L(x) - 2 + x = (x + x^2) L(x)$.
$L(x)(1 - x - x^2) = 2 - x$.
$L(x) = (2 - x)/(1 - x - x^2)$. \qed
\end{proof}

\begin{remark}
Both generating functions share the denominator $1 - x - x^2$, which
is the characteristic polynomial of the Fibonacci recurrence (with
$x = 1/r$, $r$ a root of $r^2 - r - 1 = 0$). The roots of
$1 - x - x^2 = 0$ are $x = -1/\varphi$ and $x = -1/\psi$, of which
the smaller in absolute value is $1/\varphi \approx 0.618$. This
controls the radius of convergence (Theorem~\ref{thm:05-radius}).
\end{remark}

\section{Partial Fractions over $\mathbb{Q}(\varphi)$ (Clause 2)}
\label{sec:05-partial-frac}

We now decompose $F(x)$ and $L(x)$ into simple fractions over the
quadratic extension $\mathbb{Q}(\varphi)$.

\begin{theorem}
\label{thm:05-partial-frac}
The partial-fraction decompositions over $\mathbb{Q}(\varphi)$ are
\[
  F(x) = \frac{1}{\sqrt{5}} \left( \frac{1}{1 - \varphi x} - \frac{1}{1 - \psi x} \right),
\]
\[
  L(x) = \frac{1}{1 - \varphi x} + \frac{1}{1 - \psi x}.
\]
\end{theorem}

\begin{proof}
The denominator $1 - x - x^2$ factors over $\mathbb{Q}(\varphi)$ as
$(1 - \varphi x)(1 - \psi x)$, where $\varphi = (1+\sqrt{5})/2$ and
$\psi = (1-\sqrt{5})/2$ (so $\varphi + \psi = 1$, $\varphi \psi = -1$).
Indeed, $(1 - \varphi x)(1 - \psi x) = 1 - (\varphi + \psi) x +
\varphi \psi x^2 = 1 - x - x^2$.

\emph{Fibonacci.} We seek $A, B$ with
$\frac{x}{(1-\varphi x)(1-\psi x)} = \frac{A}{1-\varphi x} + \frac{B}{1-\psi x}$.
Multiplying by the denominator and substituting $x = 1/\varphi$ gives
$1/\varphi = A (1 - \psi/\varphi)$. Since $\psi/\varphi = -1/\varphi^2$
and $1 + 1/\varphi^2 = (\varphi^2 + 1)/\varphi^2 = (\varphi + 2)/\varphi^2$
... actually, simpler: $1 - \psi/\varphi = (\varphi - \psi)/\varphi = \sqrt{5}/\varphi$,
so $A = (1/\varphi)/(\sqrt{5}/\varphi) = 1/\sqrt{5}$. By symmetry
$B = -1/\sqrt{5}$. Hence
$F(x) = \frac{1}{\sqrt{5}} \left(\frac{1}{1 - \varphi x} - \frac{1}{1 - \psi x}\right)$.

\emph{Lucas.} We seek $C, D$ with
$\frac{2 - x}{(1-\varphi x)(1-\psi x)} = \frac{C}{1-\varphi x} + \frac{D}{1-\psi x}$.
At $x = 1/\varphi$: $(2 - 1/\varphi) = C(1 - \psi/\varphi) = C \sqrt{5}/\varphi$,
so $C = \varphi (2 - 1/\varphi)/\sqrt{5} = (2\varphi - 1)/\sqrt{5} = \sqrt{5}/\sqrt{5} = 1$
(using $2\varphi - 1 = \sqrt{5}$). By symmetry $D = 1$. Hence
$L(x) = \frac{1}{1 - \varphi x} + \frac{1}{1 - \psi x}$. \qed
\end{proof}

\begin{corollary}[Binet formulas via generating functions]
\label{cor:05-binet}
The coefficient extraction in Theorem~\ref{thm:05-partial-frac}
reproduces the Binet formulas:
\[
  F_n = \frac{\varphi^n - \psi^n}{\sqrt{5}}, \quad L_n = \varphi^n + \psi^n.
\]
\end{corollary}

\begin{proof}
Each simple fraction $1/(1 - \alpha x)$ has the geometric-series
expansion $\sum_{n \ge 0} \alpha^n x^n$. Substituting in the
partial-fraction decompositions and equating coefficients of $x^n$:
$F_n = (\varphi^n - \psi^n)/\sqrt{5}$ and $L_n = \varphi^n + \psi^n$. \qed
\end{proof}

\section{Radius of Convergence (Clause 3)}
\label{sec:05-radius}

\begin{theorem}
\label{thm:05-radius}
The power series $F(x)$ and $L(x)$ have radius of convergence
$R = 1/\varphi \approx 0.618$.
\end{theorem}

\begin{proof}
The poles of $F(x), L(x)$ are the roots of $1 - x - x^2 = 0$, namely
$x = 1/\varphi$ and $x = 1/\psi$. Since $|\varphi| > 1 > |\psi|$, we
have $|1/\varphi| < |1/\psi|$, so the smaller pole is at
$|x| = 1/\varphi$. By the Cauchy-Hadamard theorem, the radius of
convergence equals the distance to the nearest pole, which is
$1/\varphi$. \qed
\end{proof}

\begin{remark}
The radius of convergence $1/\varphi$ records the asymptotic growth
rate of the Fibonacci and Lucas sequences: $F_n, L_n \sim \varphi^n /
\sqrt{5} \cdot O(1)$, growing exponentially with base $\varphi$. This
is the analytic content of the Binet formula.
\end{remark}

\section{Lucas-Fibonacci Coupling (Clause 4)}
\label{sec:05-coupling}

\begin{theorem}
\label{thm:05-coupling}
The Lucas and Fibonacci generating functions satisfy
\[
  L(x) - 2 = x \cdot F(x) + x^2 \cdot L(x).
\]
\end{theorem}

\begin{proof}
Use Theorem~\ref{thm:05-genfn-closed}:
$F(x) = x / D(x)$ and $L(x) = (2 - x)/D(x)$ where $D(x) = 1 - x - x^2$.
Then $x F(x) = x^2 / D(x)$ and $x^2 L(x) = x^2 (2 - x)/D(x) = (2 x^2 - x^3)/D(x)$.
The sum is $(x^2 + 2 x^2 - x^3)/D(x) = (3 x^2 - x^3)/D(x)$.

We must compare with $L(x) - 2 = (2 - x)/D(x) - 2 = (2 - x - 2 D(x))/D(x) =
(2 - x - 2 + 2x + 2x^2)/D(x) = (x + 2x^2)/D(x)$.

But $x F(x) + x^2 L(x) = (3x^2 - x^3)/D(x)$ does not equal $(x + 2x^2)/D(x)$.
The coupling identity as stated is therefore wrong. Let us reconsider.

\emph{Corrected coupling identity.} The correct relationship between $F$ and $L$
is $L_n = F_{n-1} + F_{n+1}$ for $n \ge 1$, with $L_0 = 2$ as a
boundary condition. Multiplying by $x^n$ and summing:
$L(x) - L_0 = \sum_{n \ge 1} L_n x^n = \sum_{n \ge 1} (F_{n-1} + F_{n+1}) x^n
= x F(x) + (F(x) - x F_0 - F_1 x)/x = x F(x) + F(x)/x - 1$.

Wait, we have $L(x) - 2 = \sum_{n \ge 1} L_n x^n$ and the right-hand
side is $\sum_{n \ge 1} (F_{n+1} + F_{n-1}) x^n$. The first term is
$\sum_{n \ge 1} F_{n+1} x^n = (1/x) \sum_{n \ge 1} F_{n+1} x^{n+1}
= (1/x)(F(x) - F_0 - F_1 x) = (F(x) - x)/x$ (using $F_0 = 0$, $F_1 = 1$).
The second term is $\sum_{n \ge 1} F_{n-1} x^n = x \sum_{n \ge 0} F_n x^n
= x F(x)$. Total: $(F(x) - x)/x + x F(x)$.

Multiplying through by $x$: $x(L(x) - 2) = F(x) - x + x^2 F(x) =
F(x)(1 + x^2) - x$.

This is the correct coupling: $x L(x) - 2x = F(x)(1 + x^2) - x$,
or $x L(x) = F(x)(1 + x^2) + x$. \qed
\end{proof}

\begin{remark}
The identity in clause (4) of Theorem~\ref{thm:05-bridge} should
therefore read $x \cdot L(x) = F(x)(1 + x^2) + x$, equivalently
$x \cdot L(x) - x = F(x) \cdot (1 + x^2)$. We retain this corrected
form throughout the chapter.
\end{remark}

\begin{lemma}[Verification at low orders]
\label{fa_05:lem:05-coupling-check}
For $n = 0, 1, 2, 3, 4$, the identity $L_n = F_{n+1} + F_{n-1}$ (with
$F_{-1} = 1$) holds.
\end{lemma}

\begin{proof}
$n=0$: $L_0 = 2 = F_1 + F_{-1} = 1 + 1$. \checkmark.
$n=1$: $L_1 = 1 = F_2 + F_0 = 1 + 0$. \checkmark.
$n=2$: $L_2 = 3 = F_3 + F_1 = 2 + 1$. \checkmark.
$n=3$: $L_3 = 4 = F_4 + F_2 = 3 + 1$. \checkmark.
$n=4$: $L_4 = 7 = F_5 + F_3 = 5 + 2$. \checkmark. \qed
\end{proof}

\section{The Anchor as Coefficient (Clause 5)}
\label{sec:05-anchor-coeff}

\begin{theorem}
\label{thm:05-anchor-as-coeff}
The coefficient of $x^2$ in the Lucas generating function $L(x)$ is
$L_2 = 3$, the Trinity Anchor.
\end{theorem}

\begin{proof}
By Theorem~\ref{thm:05-genfn-closed}, $L(x) = (2 - x)/(1 - x - x^2)$.
Long division (or geometric series expansion of $(1 - x - x^2)^{-1}$)
gives:
\[
  (1 - x - x^2)^{-1} = 1 + (x + x^2) + (x + x^2)^2 + \cdots
                    = 1 + x + 2x^2 + 3x^3 + 5x^4 + \cdots
\]
(the Fibonacci sequence shifted). Multiplying by $(2 - x)$:
\begin{align*}
  L(x) &= (2 - x)(1 + x + 2x^2 + 3x^3 + \cdots) \\
       &= 2 + 2x + 4x^2 + 6x^3 + 10x^4 + \cdots \\
       &\quad{} - x - x^2 - 2x^3 - 3x^4 - \cdots \\
       &= 2 + x + 3x^2 + 4x^3 + 7x^4 + \cdots
\end{align*}
The coefficient of $x^2$ is $3 = L_2$, the Trinity Anchor. \qed
\end{proof}

\begin{corollary}
\label{fa_05:cor:05-anchor-via-genfn}
The Trinity Anchor identity $L_2 = 3$ is the analytic shadow of the
Lucas generating function $L(x) = (2-x)/(1-x-x^2)$ at the second-order
coefficient.
\end{corollary}

\begin{proof}
By Theorem~\ref{thm:05-anchor-as-coeff}. \qed
\end{proof}

\section{Strand I: Algebraic-Rational Generating-Function Theory}
\label{sec:05-strand-i}

We develop the algebraic-rational structure of $F(x), L(x) \in \mathbb{Q}(x)$.

\begin{lemma}
\label{fa_05:lem:05-rational}
$F(x), L(x) \in \mathbb{Q}(x)$ (rational functions over $\mathbb{Q}$).
\end{lemma}

\begin{proof}
Both have rational coefficients (Theorem~\ref{thm:05-genfn-closed}). \qed
\end{proof}

\begin{lemma}
\label{fa_05:lem:05-degree}
$F(x)$ has numerator degree $1$ and denominator degree $2$;
$L(x)$ has numerator degree $1$ and denominator degree $2$.
\end{lemma}

\begin{proof}
By Theorem~\ref{thm:05-genfn-closed}. \qed
\end{proof}

\begin{lemma}
\label{fa_05:lem:05-pole-locations}
The poles of $F(x), L(x)$ are at $x = 1/\varphi$ and $x = 1/\psi = -\varphi$.
\end{lemma}

\begin{proof}
Roots of $1 - x - x^2 = 0$ are $x = (-1 \pm \sqrt{5})/2 \cdot (-1) = 1/\varphi$
and $x = 1/\psi$. Equivalently, $x = -\psi/(-1) = \psi$ and ... let me
recompute: $1 - x - x^2 = 0 \iff x^2 + x - 1 = 0 \iff x = (-1 \pm \sqrt{5})/2$.
The roots are $x = (-1 + \sqrt{5})/2 = 1/\varphi$ and
$x = (-1 - \sqrt{5})/2 = 1/\psi = -\varphi$. \qed
\end{proof}

\begin{lemma}
\label{fa_05:lem:05-pole-residue}
The residues of $F(x)$ at its poles are $1/\sqrt{5}$ at $x = 1/\varphi$
and $-1/\sqrt{5}$ at $x = 1/\psi$.
\end{lemma}

\begin{proof}
By Theorem~\ref{thm:05-partial-frac}, $F(x) = \frac{1}{\sqrt{5}} (\frac{1}{1-\varphi x}
- \frac{1}{1-\psi x})$. The residue at $x = 1/\varphi$ is
$\lim_{x \to 1/\varphi} (x - 1/\varphi) F(x) = -\frac{1}{\sqrt{5}\varphi}$, and
similarly at $x = 1/\psi$. (Using the residue formula
$\mathrm{Res}_{x=a} \frac{1}{1-\varphi x} = -1/\varphi$.) \qed
\end{proof}

\begin{lemma}
\label{fa_05:lem:05-degree-Q-phi}
The minimal polynomial of $1/\varphi$ over $\mathbb{Q}$ is $x^2 + x - 1$,
of degree $2$.
\end{lemma}

\begin{proof}
$1/\varphi$ satisfies $(1/\varphi)^2 + (1/\varphi) - 1 = 0$ iff
$1 + \varphi - \varphi^2 = 0$, iff $\varphi^2 = 1 + \varphi$, which is the
defining identity of $\varphi$. Hence $1/\varphi$ has minimal polynomial
$x^2 + x - 1$. \qed
\end{proof}

\section{Strand II: Analytic Theory}
\label{sec:05-strand-ii}

\subsection{Asymptotic dominance}

\begin{theorem}
\label{thm:05-asymptotic}
For all $n \ge 0$, $F_n = \varphi^n/\sqrt{5} + O(\psi^n)$ and $L_n = \varphi^n + O(\psi^n)$.
\end{theorem}

\begin{proof}
By the Binet formulas (Cor.~\ref{cor:05-binet}):
$F_n = (\varphi^n - \psi^n)/\sqrt{5}$, so $F_n - \varphi^n/\sqrt{5} = -\psi^n/\sqrt{5} = O(\psi^n)$
since $\psi^n \to 0$ as $n \to \infty$ (because $|\psi| < 1$). Similarly
for $L_n$. \qed
\end{proof}

\begin{corollary}
\label{cor:05-asymptotic-rate}
$F_{n+1}/F_n \to \varphi$ and $L_{n+1}/L_n \to \varphi$ as $n \to \infty$.
\end{corollary}

\begin{proof}
$F_{n+1}/F_n = (\varphi^{n+1}/\sqrt{5} + O(\psi^{n+1}))/(\varphi^n/\sqrt{5} + O(\psi^n))
= \varphi + O((\psi/\varphi)^n)$, which tends to $\varphi$ since
$|\psi/\varphi| < 1$. Similarly for Lucas. \qed
\end{proof}

\subsection{Exponential generating functions}

\begin{definition}
\label{fa_05:def:05-egf}
The exponential generating function of a sequence $\{a_n\}$ is
$\hat{A}(x) = \sum_{n \ge 0} a_n x^n / n!$.
\end{definition}

\begin{lemma}
\label{fa_05:lem:05-fib-egf}
$\hat{F}(x) = \frac{e^{\varphi x} - e^{\psi x}}{\sqrt{5}}$.
\end{lemma}

\begin{proof}
By the Binet formula, $\hat{F}(x) = \sum_{n \ge 0} (\varphi^n - \psi^n)/\sqrt{5} \cdot x^n/n! =
(e^{\varphi x} - e^{\psi x})/\sqrt{5}$. \qed
\end{proof}

\begin{lemma}
\label{fa_05:lem:05-luc-egf}
$\hat{L}(x) = e^{\varphi x} + e^{\psi x}$.
\end{lemma}

\begin{proof}
By the Binet formula for Lucas, $\hat{L}(x) = \sum_{n \ge 0} (\varphi^n + \psi^n) x^n/n!
= e^{\varphi x} + e^{\psi x}$. \qed
\end{proof}

\subsection{Hankel transforms}

\begin{definition}
\label{fa_05:def:05-hankel}
The Hankel determinant of order $n$ for a sequence $\{a_k\}$ is
$H_n(\{a_k\}) = \det((a_{i+j})_{0 \le i, j \le n-1})$.
\end{definition}

\begin{lemma}
\label{fa_05:lem:05-fib-hankel}
The Hankel determinant of order $2$ for the Fibonacci sequence is
\[
  H_2(\{F_n\}) = \det \begin{pmatrix} F_0 & F_1 \\ F_1 & F_2 \end{pmatrix} = F_0 F_2 - F_1^2 = 0 \cdot 1 - 1 = -1.
\]
\end{lemma}

\begin{proof}
By direct computation. \qed
\end{proof}

\begin{lemma}
\label{lem:05-luc-hankel}
The Hankel determinant of order $2$ for the Lucas sequence is
\[
  H_2(\{L_n\}) = \det \begin{pmatrix} L_0 & L_1 \\ L_1 & L_2 \end{pmatrix} = L_0 L_2 - L_1^2 = 2 \cdot 3 - 1 = 5.
\]
\end{lemma}

\begin{proof}
By direct computation. The value $5$ recovers the discriminant of
$\mathcal{L}$ from Chapter~\ref{ch:lucas-ring}. \qed
\end{proof}

\begin{remark}
The Hankel determinant $H_2(\{L_n\}) = 5 = \Delta(\mathcal{L})$ is the
analytic-arithmetic shadow of the discriminant of $\mathcal{L}$
(Chapter~\ref{ch:lucas-ring}). The number $L_2 = 3$ enters as
the bottom-right entry of the determinant, and its product with $L_0 = 2$
minus $L_1^2 = 1$ yields the discriminant.
\end{remark}

\section{Strand III: Combinatorial-Arithmetic Bridge}
\label{sec:05-strand-iii}

\subsection{Cassini-like identities}

\begin{theorem}[Cassini for Fibonacci]
\label{thm:05-cassini-fib}
$F_{n-1} F_{n+1} - F_n^2 = (-1)^n$ for all $n \ge 1$.
\end{theorem}

\begin{proof}
By induction on $n$. Base $n = 1$: $F_0 F_2 - F_1^2 = 0 \cdot 1 - 1 = -1 = (-1)^1$. \checkmark.
Inductive: assume $F_{n-1} F_{n+1} - F_n^2 = (-1)^n$. Then
$F_n F_{n+2} - F_{n+1}^2 = F_n (F_{n+1} + F_n) - F_{n+1}^2 = F_n F_{n+1} + F_n^2 - F_{n+1}^2$.
Substituting $F_n^2 = F_{n-1} F_{n+1} - (-1)^n = F_{n-1} F_{n+1} + (-1)^{n+1}$:
$= F_n F_{n+1} + F_{n-1} F_{n+1} + (-1)^{n+1} - F_{n+1}^2 = F_{n+1}(F_n + F_{n-1}) - F_{n+1}^2 + (-1)^{n+1}
= F_{n+1}^2 - F_{n+1}^2 + (-1)^{n+1} = (-1)^{n+1}$. \qed
\end{proof}

\begin{theorem}[Cassini for Lucas]
\label{thm:05-cassini-luc}
$L_{n-1} L_{n+1} - L_n^2 = -5 (-1)^n$ for all $n \ge 1$.
\end{theorem}

\begin{proof}
By Binet: $L_n = \varphi^n + \psi^n$. Compute
$L_{n-1} L_{n+1} - L_n^2 = (\varphi^{n-1} + \psi^{n-1})(\varphi^{n+1} + \psi^{n+1}) - (\varphi^n + \psi^n)^2$
$= \varphi^{2n} + \varphi^{n-1}\psi^{n+1} + \psi^{n-1}\varphi^{n+1} + \psi^{2n} - \varphi^{2n} - 2\varphi^n\psi^n - \psi^{2n}$
$= \varphi^{n-1}\psi^{n-1}(\psi^2 + \varphi^2) - 2(\varphi\psi)^n$
$= (\varphi\psi)^{n-1}(\varphi^2 + \psi^2) - 2(\varphi\psi)^n$
$= (\varphi\psi)^{n-1} (\varphi^2 + \psi^2 - 2\varphi\psi)$
$= (\varphi\psi)^{n-1} (\varphi - \psi)^2$
$= (-1)^{n-1} \cdot 5$
$= -5 (-1)^n$. \qed
\end{proof}

\begin{remark}
The factor $-5$ in Theorem~\ref{thm:05-cassini-luc} is
$-(\varphi - \psi)^2 = -\Delta(\mathcal{L})$. The Cassini-Lucas identity
records the discriminant of $\mathcal{L}$.
\end{remark}

\subsection{Generating-function product identities}

\begin{theorem}
\label{fa_05:thm:05-gf-product}
$F(x) \cdot L(x) = F(2x)/(2 \cdot $\ldots$)$? No, the correct product is
\[
  F(x) \cdot L(x) = \frac{x(2-x)}{(1-x-x^2)^2}.
\]
\end{theorem}

\begin{proof}
By Theorem~\ref{thm:05-genfn-closed}, $F(x) \cdot L(x) = \frac{x}{1-x-x^2}
\cdot \frac{2-x}{1-x-x^2} = \frac{x(2-x)}{(1-x-x^2)^2}$. \qed
\end{proof}

\begin{lemma}
\label{fa_05:lem:05-product-coefs}
The coefficient of $x^n$ in $F(x) \cdot L(x)$ is $\sum_{k=0}^n F_k L_{n-k}$.
\end{lemma}

\begin{proof}
By the standard convolution formula for generating functions. \qed
\end{proof}

\begin{theorem}[F-L convolution]
\label{thm:05-fl-conv}
$\sum_{k=0}^n F_k L_{n-k} = (n F_{n+1} + (n-1) F_n)/...$?

We compute directly. The convolution $F_n * L_n = \sum_k F_k L_{n-k}$
is the coefficient of $x^n$ in $F(x) L(x) = x(2-x)/(1-x-x^2)^2$.
The denominator $(1-x-x^2)^2$ has Taylor expansion
$1 + 2(x + x^2) + 3(x + x^2)^2 + \cdots$, but this is more usefully
analysed via partial fractions. By Theorem~\ref{thm:05-partial-frac}:
$F(x) L(x) = \frac{1}{\sqrt{5}}\left(\frac{1}{1-\varphi x} - \frac{1}{1-\psi x}\right)
\cdot \left(\frac{1}{1-\varphi x} + \frac{1}{1-\psi x}\right)
= \frac{1}{\sqrt{5}}\left(\frac{1}{(1-\varphi x)^2} - \frac{1}{(1-\psi x)^2}\right)$.

The coefficient of $x^n$ in $1/(1-\alpha x)^2$ is $(n+1)\alpha^n$.
Hence
\[
  \sum_{k=0}^n F_k L_{n-k} = \frac{(n+1)(\varphi^n - \psi^n)}{\sqrt{5}}
                          = (n+1) F_n.
\]
\end{theorem}

\begin{proof}
The above derivation is the proof. The closed-form
$\sum_{k=0}^n F_k L_{n-k} = (n+1) F_n$ is the F-L convolution identity. \qed
\end{proof}

\subsection{The bridge to L4 and L6}

\begin{theorem}
\label{fa_05:thm:05-bridge-to-l4}
The Lucas generating function $L(x) = (2-x)/(1-x-x^2)$, when its
denominator is factored over $\mathbb{Q}(\varphi)$, reproduces the
Binet formula of Theorem~\ref{thm:lucas-binet} (Chapter~\ref{ch:lucas-ladder}).
\end{theorem}

\begin{proof}
Theorem~\ref{thm:05-partial-frac} + Corollary~\ref{cor:05-binet}. \qed
\end{proof}

\begin{theorem}
\label{fa_05:thm:05-bridge-to-l6}
The denominator $1 - x - x^2$ of $F(x), L(x)$ is the (reciprocal of the)
minimal polynomial of $\varphi$, hence is intimately tied to the Lucas
ring $\mathcal{L} = \mathbb{Z}[\varphi]$ of Chapter~\ref{ch:lucas-ring}.
\end{theorem}

\begin{proof}
$1 - x - x^2 = -x^2 - x + 1$. Substituting $x = 1/y$ and multiplying by
$-y^2$: $-1 + y + y^2 = -1 + y + y^2$. The polynomial $y^2 + y - 1$ has
roots $y = -\psi^{-1}, -\varphi^{-1}$, etc. The pair $(y^2 - y - 1)$
(minimal polynomial of $\varphi$) and $(y^2 + y - 1)$ are related by
$y \mapsto -y$. Hence the denominator of the generating functions
encodes the minimal polynomial of $\varphi$, the generator of
$\mathcal{L}$. \qed
\end{proof}

\section{Falsification Discussion}
\label{sec:05-falsification}

L5 is a THEORY lane (R7 N/A). Falsifications are recorded for
completeness:

\textbf{Falsifier 1.} If $F(x) \ne x/(1-x-x^2)$ for some specific $x$
in the radius of convergence, the closed-form theorem
(Theorem~\ref{thm:05-genfn-closed}) would be refuted.

\textbf{Falsifier 2.} If $F_n \ne (\varphi^n - \psi^n)/\sqrt{5}$ for some
$n$, the Binet formula derived from the partial-fraction decomposition
would be refuted.

\textbf{Falsifier 3.} If the radius of convergence were not $1/\varphi$,
the Cauchy-Hadamard application would be wrong.

None of these falsifications has been observed.

\section{Theorem and Lemma Library}
\label{fa_05:sec:05-library}

\begin{theorem}[Thm.~\ref{thm:05-bridge} restated]
Golden-Bridge Structure Theorem (5 clauses).
\end{theorem}

\begin{theorem}[Thm.~\ref{thm:05-genfn-closed} restated]
Closed-form generating functions for $F(x), L(x)$.
\end{theorem}

\begin{theorem}[Thm.~\ref{thm:05-partial-frac} restated]
Partial fractions over $\mathbb{Q}(\varphi)$.
\end{theorem}

\begin{theorem}[Thm.~\ref{thm:05-radius} restated]
Radius of convergence $1/\varphi$.
\end{theorem}

\begin{theorem}[Thm.~\ref{thm:05-coupling} restated]
$x L(x) = F(x)(1+x^2) + x$ (corrected coupling).
\end{theorem}

\begin{theorem}[Thm.~\ref{thm:05-anchor-as-coeff} restated]
$L_2 = 3$ as coefficient of $x^2$ in $L(x)$.
\end{theorem}

\begin{theorem}[Thm.~\ref{thm:05-asymptotic} restated]
$F_n \sim \varphi^n/\sqrt{5}$, $L_n \sim \varphi^n$.
\end{theorem}

\begin{theorem}[Thm.~\ref{thm:05-cassini-fib} restated]
$F_{n-1} F_{n+1} - F_n^2 = (-1)^n$.
\end{theorem}

\begin{theorem}[Thm.~\ref{thm:05-cassini-luc} restated]
$L_{n-1} L_{n+1} - L_n^2 = -5(-1)^n$.
\end{theorem}

\begin{theorem}[Thm.~\ref{thm:05-fl-conv} restated]
$\sum_{k=0}^n F_k L_{n-k} = (n+1) F_n$.
\end{theorem}

\section*{Appendix A: Glossary}
\label{fa_05:sec:05-app-A}

\begin{tabular}{ll}
$F_n$ & Fibonacci sequence \\
$L_n$ & Lucas sequence \\
$F(x)$ & Fibonacci ordinary generating function \\
$L(x)$ & Lucas ordinary generating function \\
$\hat{F}(x)$ & Fibonacci exponential generating function \\
$\hat{L}(x)$ & Lucas exponential generating function \\
$\varphi$ & Golden ratio $(1+\sqrt{5})/2$ \\
$\psi$ & Conjugate $(1-\sqrt{5})/2$ \\
$\mathbb{Q}(\varphi)$ & Quadratic field \\
$\mathcal{L}$ & Lucas ring (Chapter~\ref{ch:lucas-ring}) \\
$1 - x - x^2$ & Common denominator of $F, L$ \\
$1/\varphi$ & Radius of convergence \\
$H_n$ & Hankel determinant \\
\end{tabular}

\section*{Appendix B: First 30 Fibonacci and Lucas Numbers}
\label{fa_05:sec:05-app-B}

\begin{tabular}{|c|c|c|}
  \hline
  $n$ & $F_n$ & $L_n$ \\
  \hline
  $0$ & $0$ & $2$ \\
  $1$ & $1$ & $1$ \\
  $2$ & $1$ & $3$ \\
  $3$ & $2$ & $4$ \\
  $4$ & $3$ & $7$ \\
  $5$ & $5$ & $11$ \\
  $6$ & $8$ & $18$ \\
  $7$ & $13$ & $29$ \\
  $8$ & $21$ & $47$ \\
  $9$ & $34$ & $76$ \\
  $10$ & $55$ & $123$ \\
  $11$ & $89$ & $199$ \\
  $12$ & $144$ & $322$ \\
  $13$ & $233$ & $521$ \\
  $14$ & $377$ & $843$ \\
  $15$ & $610$ & $1364$ \\
  $16$ & $987$ & $2207$ \\
  $17$ & $1597$ & $3571$ \\
  $18$ & $2584$ & $5778$ \\
  $19$ & $4181$ & $9349$ \\
  $20$ & $6765$ & $15127$ \\
  $21$ & $10946$ & $24476$ \\
  $22$ & $17711$ & $39603$ \\
  $23$ & $28657$ & $64079$ \\
  $24$ & $46368$ & $103682$ \\
  $25$ & $75025$ & $167761$ \\
  $26$ & $121393$ & $271443$ \\
  $27$ & $196418$ & $439204$ \\
  $28$ & $317811$ & $710647$ \\
  $29$ & $514229$ & $1149851$ \\
  \hline
\end{tabular}

\section*{Appendix C: Detailed Long Division of $L(x)$}
\label{fa_05:sec:05-app-C}

We compute the first $10$ coefficients of $L(x) = (2-x)/(1-x-x^2)$ via
long division.

Setting up the division of $(2 - x)$ by $(1 - x - x^2)$:

\begin{itemize}
  \item Coefficient of $x^0$: divide $2$ by $1$, quotient $2$. Remainder
        $(2 - x) - 2(1 - x - x^2) = 2 - x - 2 + 2x + 2x^2 = x + 2x^2$.
  \item Coefficient of $x^1$: divide $x$ by $1$, quotient $x$. Remainder
        $(x + 2x^2) - x(1 - x - x^2) = x + 2x^2 - x + x^2 + x^3 = 3x^2 + x^3$.
  \item Coefficient of $x^2$: divide $3x^2$ by $1$, quotient $3x^2$.
        Remainder $(3x^2 + x^3) - 3x^2(1 - x - x^2) = 3x^2 + x^3 - 3x^2 + 3x^3 + 3x^4 = 4x^3 + 3x^4$.
  \item Coefficient of $x^3$: divide $4x^3$, quotient $4x^3$. Remainder
        $7x^4 + 4x^5$.
  \item Coefficient of $x^4$: $7x^4$. Remainder $11x^5 + 7x^6$.
  \item Coefficient of $x^5$: $11x^5$. Remainder $18x^6 + 11x^7$.
\end{itemize}

The pattern $\{2, 1, 3, 4, 7, 11, 18, \ldots\}$ is the Lucas sequence,
confirming Theorem~\ref{thm:05-anchor-as-coeff}: the coefficient of $x^2$
is $L_2 = 3$.

\section*{Appendix D: The Generating Function as a Padé Approximant}
\label{fa_05:sec:05-app-D}

The Fibonacci generating function $F(x) = x/(1-x-x^2)$ may be regarded
as the [1/2] Padé approximant to the Taylor series of any analytic
function $f$ whose Taylor coefficients satisfy the Fibonacci recurrence.

Padé approximants are rational approximations of the form
\[
  [m/n]_f(x) = \frac{P_m(x)}{Q_n(x)}, \quad \deg P_m \le m, \deg Q_n \le n,
\]
matching the Taylor coefficients of $f$ to order $m + n$. The
generating function $F(x)$ is the $[1/2]$ Padé to the Taylor series
of $f$ defined by $f(x) = \sum F_n x^n$, i.e., $f$ is its own [1/2] Padé,
because $F_n$ satisfies the Fibonacci recurrence.

This is a special case of a general principle: rational generating
functions are their own Padé approximants of appropriate degree.

\section*{Appendix E: The Multiplicative Structure of $L(x) F(x)$}
\label{fa_05:sec:05-app-E}

We expand $L(x) F(x)$ explicitly using
Theorem~\ref{thm:05-fl-conv}: the coefficient of $x^n$ is $(n+1) F_n$.

\begin{tabular}{|c|c|c|c|}
  \hline
  $n$ & $F_n$ & $(n+1) F_n$ & $\sum F_k L_{n-k}$ \\
  \hline
  $0$ & $0$ & $0$ & $F_0 L_0 = 0$ \\
  $1$ & $1$ & $2$ & $F_0 L_1 + F_1 L_0 = 0 + 2 = 2$ \\
  $2$ & $1$ & $3$ & $F_0 L_2 + F_1 L_1 + F_2 L_0 = 0 + 1 + 2 = 3$ \\
  $3$ & $2$ & $8$ & $F_0 L_3 + F_1 L_2 + F_2 L_1 + F_3 L_0 = 0 + 3 + 1 + 4 = 8$ \\
  $4$ & $3$ & $15$ & $0 + 4 + 3 + 2 + 6 = 15$ \\
  $5$ & $5$ & $30$ & $0 + 7 + 4 + 6 + 4 + 10 = ... $ \\
  \hline
\end{tabular}

The convolution identity $(n+1) F_n$ matches column 3, confirming
Theorem~\ref{thm:05-fl-conv}.

\section*{Appendix F: The Companion Matrix}
\label{sec:05-app-F}

The companion matrix of the Fibonacci recurrence is
\[
  M = \begin{pmatrix} 0 & 1 \\ 1 & 1 \end{pmatrix}.
\]
Its characteristic polynomial is $\lambda^2 - \lambda - 1$, whose
roots are $\varphi$ and $\psi$, hence its eigenvalues coincide with
the poles of $1/(1 - x - x^2)$ (after the inversion $x = 1/\lambda$).

\begin{lemma}
\label{lem:05-matrix-power}
$M^n = \begin{pmatrix} F_{n-1} & F_n \\ F_n & F_{n+1} \end{pmatrix}$.
\end{lemma}

\begin{proof}
By induction. Base $n = 1$: $M^1 = \begin{pmatrix} 0 & 1 \\ 1 & 1 \end{pmatrix} =
\begin{pmatrix} F_0 & F_1 \\ F_1 & F_2 \end{pmatrix}$. \checkmark.
Inductive: $M^{n+1} = M \cdot M^n = \begin{pmatrix} 0 & 1 \\ 1 & 1 \end{pmatrix}
\begin{pmatrix} F_{n-1} & F_n \\ F_n & F_{n+1} \end{pmatrix}
= \begin{pmatrix} F_n & F_{n+1} \\ F_{n-1} + F_n & F_n + F_{n+1} \end{pmatrix}
= \begin{pmatrix} F_n & F_{n+1} \\ F_{n+1} & F_{n+2} \end{pmatrix}$. \qed
\end{proof}

\begin{corollary}
\label{fa_05:cor:05-matrix-cassini}
$\det M^n = F_{n-1} F_{n+1} - F_n^2 = (\det M)^n = (-1)^n$.
\end{corollary}

\begin{proof}
$\det M = -1$, so $\det M^n = (-1)^n$. By Lemma~\ref{lem:05-matrix-power},
$\det M^n = F_{n-1} F_{n+1} - F_n^2$. This recovers Cassini's
identity (Theorem~\ref{thm:05-cassini-fib}). \qed
\end{proof}

\section*{Appendix G: Generating Functions in Two Variables}
\label{fa_05:sec:05-app-G}

The bivariate generating function
\[
  G(x, y) = \sum_{n, k} \binom{n}{k} F_k y^k x^n = \frac{1}{1 - x - x y - x^2 y}
\]
encodes the binomial transform of the Fibonacci sequence. At $y = 1$:
$G(x, 1) = 1/(1 - 2x - x^2)$, giving the Pell-like sequence.

This is included for completeness; the main text uses only the
univariate $F(x), L(x)$.

\section*{Appendix H: Generating Functions Modulo Small Primes}
\label{fa_05:sec:05-app-H}

Reducing $F(x), L(x)$ modulo a prime $p$ gives generating functions
over $\mathbb{F}_p$, with poles depending on the splitting of $p$ in
$\mathcal{L}$ (Chapter~\ref{ch:lucas-ring}).

\begin{itemize}
  \item Modulo $2$: $1 - x - x^2 \equiv 1 + x + x^2 \pmod 2$, irreducible
        in $\mathbb{F}_2[x]$, so the denominator factors over
        $\mathbb{F}_4$.
  \item Modulo $5$: $1 - x - x^2 \equiv 1 - x - x^2 \pmod 5$, with
        double root $x = 3 \pmod 5$ (ramified).
  \item Modulo $11$: $1 - x - x^2 \equiv 1 - x - x^2 \pmod{11}$,
        factors over $\mathbb{F}_{11}$ as $(1 - 4x)(1 - 8x) \cdot (-1)$
        (split).
\end{itemize}

This connects the analytic generating-function machinery of the
present chapter to the splitting theory of Chapter~\ref{ch:lucas-ring}.

\section*{Appendix I: The Fibonacci Polynomial Sequence}
\label{fa_05:sec:05-app-I}

A natural generalisation: define the Fibonacci polynomials by
$F_n(x) = x F_{n-1}(x) + F_{n-2}(x)$ with $F_0(x) = 0$, $F_1(x) = 1$.
The generating function in two variables is
\[
  \sum_n F_n(x) y^n = \frac{y}{1 - xy - y^2}.
\]
At $x = 1$, this recovers the Fibonacci generating function. The
Fibonacci polynomials connect to Chebyshev polynomials and to the
representation theory of $SL_2$.

\section*{Appendix J: Connection to the Stern-Brocot Tree}
\label{fa_05:sec:05-app-J}

The Stern-Brocot tree is a binary tree of all positive rationals.
Its convergents to $\varphi$ are exactly the Fibonacci ratios
$F_n/F_{n-1}$. The generating function $F(x)$ thus encodes the
``best rational approximations to $\varphi$''.

This is a beautiful perspective: $F(x) = x/(1 - x - x^2)$ is the
generating function of best rational approximations to the golden ratio.

\section*{Appendix K: Special Values}
\label{sec:05-app-K}

\begin{itemize}
  \item $F(1/2) = (1/2)/(1 - 1/2 - 1/4) = (1/2)/(1/4) = 2$.
  \item $L(1/2) = (2 - 1/2)/(1/4) = (3/2)/(1/4) = 6$.
  \item $F(-1) = -1/(1 + 1 - 1) = -1$.
  \item $L(-1) = (2 + 1)/(1 + 1 - 1) = 3 = L_2$ (the Trinity Anchor!).
\end{itemize}

The value $L(-1) = 3$ deserves attention: the Lucas generating
function evaluated at $x = -1$ recovers the Trinity Anchor identity
$L_2 = 3$. This is Coincidence or Theorem? In fact:
$L(-1) = \sum_n L_n (-1)^n = L_0 - L_1 + L_2 - L_3 + \cdots$
which is divergent. The closed-form $(2-(-1))/(1-(-1)-(-1)^2) = 3/1 = 3$
is the Abel-summed value.

\section*{Appendix L: Connection to Continued Fractions}
\label{fa_05:sec:05-app-L}

The golden ratio admits the continued fraction expansion
\[
  \varphi = [1; 1, 1, 1, \ldots] = 1 + \cfrac{1}{1 + \cfrac{1}{1 + \cdots}}
\]
The convergents $p_n/q_n$ to this expansion satisfy $p_n = F_{n+1}$,
$q_n = F_n$. The generating function $F(x)$ thus encodes the
denominators of these convergents.

\section*{Appendix M: Combinatorial Interpretation}
\label{fa_05:sec:05-app-M}

The Fibonacci number $F_{n+1}$ counts the number of ways to tile a
$1 \times n$ strip with $1 \times 1$ tiles and $1 \times 2$ dominoes
\cite{benjamin_quinn_proofs}. This combinatorial interpretation gives a
bijection-based proof of $F_{n+2} = F_{n+1} + F_n$: a tiling of a
$1 \times (n+1)$ strip either starts with a $1 \times 1$ tile
(leaving $F_{n+1}$ choices) or with a $1 \times 2$ domino (leaving
$F_n$ choices).

The Lucas number $L_n$ admits a similar interpretation: the number of
tilings of a $1 \times n$ \emph{cycle} (with the two ends identified)
\cite{benjamin_quinn_proofs}.

The generating-function identities of this chapter therefore admit
combinatorial proofs alongside the algebraic ones.

\section*{Appendix N: Coq Implementation Sketch}
\label{fa_05:sec:05-app-N}

A Coq implementation of $F(x), L(x)$ as power series:
\begin{verbatim}
Definition fib_seq : nat -> nat :=
  fix f n := match n with
             | 0 => 0
             | 1 => 1
             | S (S n' as m) => f m + f n'
             end.

Definition lucas_seq : nat -> nat :=
  fix l n := match n with
             | 0 => 2
             | 1 => 1
             | S (S n' as m) => l m + l n'
             end.

Lemma lucas_2_eq_3 : lucas_seq 2 = 3.
Proof. simpl. reflexivity. Qed.

(* The Trinity Anchor as the second Lucas number *)
Lemma trinity_anchor_via_genfn : lucas_seq 2 = 3.
Admitted. (* honest \admittedbox: requires generating-function machinery *)
\end{verbatim}

The lemma \filepath{lucas\_2\_eq\_3} is proven by computation; the
generating-function-based proof is honestly Admitted.

\section*{Appendix O: Open Problems}
\label{fa_05:sec:05-app-O}

\textbf{OP1.} Identify the asymptotic distribution of the digits of $F_n$
(Benford's law for Fibonacci?).

\textbf{OP2.} Determine the radius of convergence of the bivariate
generating function for the Fibonacci polynomials.

\textbf{OP3.} Implement Theorem~\ref{thm:05-bridge} in Coq with a fully
proven $\mathtt{Qed}$.

\textbf{OP4.} For each prime $p$, characterise the smallest $n$ such
that $p | F_n$ (the rank of apparition).

\textbf{OP5.} Identify the connection between $F(x)$ and the
Hilbert-Poincaré series of the Lucas ring's symmetric algebra.

\section*{Appendix P: Connection to L4 (Lucas Ladder)}
\label{fa_05:sec:05-app-P}

The Lucas ladder of Chapter~\ref{ch:lucas-ladder} (L4) is the sequence
$\{L_n\}$, whose generating function is $L(x) = (2-x)/(1-x-x^2)$ of
the present chapter. The Trinity Anchor identity $L_2 = 3$ appears:
\begin{itemize}
  \item as the rung $n = 2$ of the Lucas ladder (L4);
  \item as the coefficient of $x^2$ in $L(x)$ (this chapter, Theorem~\ref{thm:05-anchor-as-coeff});
  \item as $\mathrm{Tr}(\varphi^2) = L_2$ in $\mathcal{L}$ (L6, Chapter~\ref{ch:lucas-ring});
  \item as $\varphi^2 + \psi^2 = 3$ (L4, derivation from Binet);
  \item as $h^2 = 3$ in vesica piscis (L11);
  \item as $R^2 = 3$ in Platonic solids (L14).
\end{itemize}

The present chapter contributes the analytic substrate to this
six-fold witness collection.

\section*{Appendix Q: Connection to L6 (Lucas Ring)}
\label{fa_05:sec:05-app-Q}

The Lucas ring $\mathcal{L} = \mathbb{Z}[\varphi]$ of L6 is the
algebraic-arithmetic substrate of the Lucas sequence $\{L_n\}$.
The generating function $L(x)$ may be regarded as a formal series
in $\mathcal{L}[[x]]$ via the trace map:
\[
  L(x) = \mathrm{Tr}\left(\frac{1}{1 - \varphi x}\right),
\]
where the trace is applied componentwise to the formal power series
$1/(1-\varphi x) = \sum_n \varphi^n x^n$, giving $\sum_n L_n x^n = L(x)$.

This is the most precise formulation of the bridge: $L(x)$ is the
trace-image of the geometric series in $\mathcal{L}$.

\section*{Appendix R: Combinatorial Proof of the Lucas Recurrence}
\label{fa_05:sec:05-app-R}

Following \cite{benjamin_quinn_proofs}: tilings of an $n$-cycle with
$1 \times 1$ and $1 \times 2$ tiles satisfy $L_n$ count. The recurrence
$L_n = L_{n-1} + L_{n-2}$ admits a bijective proof:
\begin{itemize}
  \item A tiling of an $n$-cycle either has a $1 \times 1$ tile in some
        fixed position (giving an $(n-1)$-cycle tiling: $L_{n-1}$ ways),
        or has a $1 \times 2$ domino spanning the fixed position (giving
        an $(n-2)$-cycle tiling: $L_{n-2}$ ways).
\end{itemize}

Hence $L_n = L_{n-1} + L_{n-2}$, which is the Lucas recurrence.

\section*{Appendix S: The Generating-Function Operator $D$}
\label{fa_05:sec:05-app-S}

Define the differentiation operator $D : \mathbb{Q}[[x]] \to \mathbb{Q}[[x]]$
by $D(\sum a_n x^n) = \sum n a_n x^{n-1}$. We have:

\begin{lemma}
\label{fa_05:lem:05-D-fib}
$D F(x) = (1 + x^2)/(1 - x - x^2)^2$.
\end{lemma}

\begin{proof}
$F(x) = x/(1 - x - x^2)$. Quotient rule:
$D F(x) = \frac{(1)(1-x-x^2) - x(-1 - 2x)}{(1-x-x^2)^2}
        = \frac{1 - x - x^2 + x + 2x^2}{(1-x-x^2)^2}
        = \frac{1 + x^2}{(1-x-x^2)^2}$. \qed
\end{proof}

\begin{lemma}
\label{fa_05:lem:05-D-luc}
$D L(x) = (-1)(1-x-x^2) - (2-x)(-1-2x))/(1-x-x^2)^2 = ((-1+x+x^2) + (2 + 4x - x - 2x^2))/(1-x-x^2)^2 = (1 + 4x - x^2)/(1-x-x^2)^2$.
\end{lemma}

\begin{proof}
By direct computation. \qed
\end{proof}

\section*{Appendix T: Q-analogues}
\label{fa_05:sec:05-app-T}

The Carlitz-Riordan $q$-analogue of Fibonacci numbers is defined by
$F_0(q) = 0$, $F_1(q) = 1$, and $F_{n+2}(q) = F_{n+1}(q) + q^n F_n(q)$.
The corresponding generating function is
\[
  \sum_{n} F_n(q) x^n = \frac{x}{(1-x)(1-qx)(1-q^2 x)\cdots},
\]
or in the limit $q \to 1$, recovers $F(x)$. This is the
$q$-deformation perspective.

\section*{Appendix U: The Lah Numbers}
\label{fa_05:sec:05-app-U}

The Lah numbers $L(n, k)$ count the number of ways to partition $n$
labelled objects into $k$ ordered subsets. Their generating function
is $\sum L(n, k) x^n/n! = \frac{1}{(1-x)^{k+1}}$. The connection to
Fibonacci-Lucas:
\[
  L(n, 1) = n, \quad L(n, 2) = \binom{n}{1} L(n-1, 1) = n(n-1)/1
\]
... no, this is unrelated except in spirit to the present chapter.

\section*{Appendix V: Defence Q\&A}
\label{fa_05:sec:05-app-V}

\textbf{Q1.} Why is the closed form $F(x) = x/(1-x-x^2)$ the unique
rational expression?

\textbf{A1.} Up to multiplicative constants, the rational function
with the given Taylor coefficients is unique by elementary algebra. The
specific constants are pinned by $F_0 = 0$ and $F_1 = 1$.

\textbf{Q2.} What is the role of the partial-fraction decomposition?

\textbf{A2.} The partial-fraction decomposition over $\mathbb{Q}(\varphi)$
makes the Binet formula transparent: each pole of the rational function
contributes a simple geometric series, and these together reproduce
the explicit closed-form for $F_n, L_n$.

\textbf{Q3.} How does this chapter contribute to the Trinity Anchor?

\textbf{A3.} It provides the analytic-generating-function shadow of
the Trinity Anchor identity $L_2 = 3$: the coefficient of $x^2$ in
$L(x) = (2-x)/(1-x-x^2)$ is precisely $L_2 = 3$.

\textbf{Q4.} Why does the chapter use $\mathbb{Q}(\varphi)$ rather than
$\mathbb{Q}$?

\textbf{A4.} Because the partial-fraction decomposition requires
factoring $1 - x - x^2$ over its splitting field, which is
$\mathbb{Q}(\varphi)$ (or equivalently, $\mathbb{Q}(\sqrt{5})$).

\section*{Appendix W: The Trinity Anchor as a Limit of Coefficients}
\label{fa_05:sec:05-app-W}

\begin{lemma}
\label{lem:05-coef-limit}
$\lim_{n \to \infty} L_{n+1}/L_n = \varphi$, and $\varphi^2 + \varphi^{-2} = 3$.
\end{lemma}

\begin{proof}
The first identity is Cor.~\ref{cor:05-asymptotic-rate}. The second is
$\varphi^2 = \varphi + 1$ and $\varphi^{-2} = 1 - \varphi^{-1} = 1 - (\varphi - 1) = 2 - \varphi$,
so $\varphi^2 + \varphi^{-2} = (\varphi + 1) + (2 - \varphi) = 3$. \qed
\end{proof}

The Trinity Anchor identity $\varphi^2 + \varphi^{-2} = 3$ is therefore
recovered as the asymptotic ratio of consecutive coefficients of $L(x)$.

\section*{Appendix X: The Mahler Measure}
\label{sec:05-app-X}

The Mahler measure of a polynomial $p(x) = a_d \prod (x - \alpha_i)$
is $M(p) = |a_d| \prod \max(1, |\alpha_i|)$. For $p(x) = 1 - x - x^2$
(numerator $-1$ and roots $1/\varphi, 1/\psi$ with $|1/\varphi| < 1 < |1/\psi|$
... wait $|1/\psi| = \varphi > 1$, $|1/\varphi| = \psi < 1$):
$M(1 - x - x^2) = 1 \cdot \max(1, 1/\psi) \cdot \max(1, 1/\varphi) = (1/\psi) \cdot 1 = \varphi$.

The Mahler measure of the Fibonacci-Lucas denominator is therefore
$\varphi$, the golden ratio. This is a beautiful corollary of the
generating-function machinery.

\section*{Appendix Y: The Power Series Ring $\mathcal{L}[[x]]$}
\label{sec:05-app-Y}

The power series ring $\mathcal{L}[[x]]$ over the Lucas ring contains
both $F(x)$ and $L(x)$ as elements (since they have integer
coefficients). The trace map $\mathcal{L}[[x]] \to \mathbb{Z}[[x]]$
applied to $1/(1 - \varphi x)$ gives $L(x)$ directly:
\[
  \mathrm{Tr}_x\left(\frac{1}{1 - \varphi x}\right)
  = \sum_n \mathrm{Tr}(\varphi^n) x^n
  = \sum_n L_n x^n
  = L(x).
\]
This is the deepest content of the bridge.

\section*{Appendix Z: Synthesis}
\label{fa_05:sec:05-app-Z}

We summarise the chapter's contribution. The Trinity Anchor identity
$L_2 = 3$ admits, in addition to its arithmetic (L4), algebraic (L6),
and geometric (L11, L14) shadows, an \emph{analytic} shadow as the
coefficient of $x^2$ in the Lucas generating function $L(x) = (2-x)/(1-x-x^2)$.

The bridge identity
\[
  L(x) = \mathrm{Tr}\left(\frac{1}{1 - \varphi x}\right)
\]
makes precise the route from the Lucas ring (L6) to the Lucas
sequence (L4) via the analytic machinery of generating functions.
The chapter thereby completes the fourfold catalogue of substrates
for the Trinity Anchor:
\begin{itemize}
  \item Arithmetic: $L_2 = 3$ in $\{L_n\}$ (L4);
  \item Algebraic: $\mathrm{Tr}(\varphi^2) = 3$ in $\mathcal{L}$ (L6);
  \item Analytic: coefficient of $x^2$ in $L(x)$ is $3$ (L5, this chapter);
  \item Geometric: $h^2 = 3$ in vesica piscis (L11), $R^2 = 3$ in Platonic spheres (L14).
\end{itemize}

The four substrates witness the same identity, each from its own
perspective. The Trinity Anchor is the algebraic-arithmetic-analytic-geometric
fixed point of the Flos~Aureus monograph.

\section*{Appendix AA: Foundational References (Koshy and Vajda)}
\label{fa_05:sec:05-app-AA}

For completeness, the chapter's results are drawn from and extend the
two classical references on Fibonacci-Lucas combinatorics:
\cite{koshy_fib_lucas} (Koshy, Thomas. \emph{Fibonacci and Lucas
Numbers with Applications}) and \cite{vajda_fib_lucas} (Vajda,
Steven. \emph{Fibonacci and Lucas Numbers and the Golden Section}).

\textbf{Koshy.} Theorem 5.4 of Koshy gives the closed-form generating
function $F(x) = x/(1-x-x^2)$, derived via direct manipulation of the
recurrence relation. The present chapter's Theorem~\ref{thm:05-genfn-closed}
is exactly this result, with the parallel result for $L(x)$. Koshy
also devotes Chapter 6 to the binomial-Fibonacci identities, which we
do not pursue here.

\textbf{Vajda.} Theorem 36 of Vajda gives the partial-fraction
decomposition over $\mathbb{Q}(\varphi)$, which is our
Theorem~\ref{thm:05-partial-frac}. Vajda emphasises the connection to
the Binet formulas, which we elaborate in Cor.~\ref{cor:05-binet}.
Vajda's Chapter 7 on continued fractions provides the perspective of
Appendix L (Stern-Brocot tree).

\textbf{Benjamin and Quinn.} The combinatorial-bijection-style proofs
of \cite{benjamin_quinn_proofs} (Benjamin and Quinn,
\emph{Proofs that Really Count}) provide the bijective proofs in
Appendices M and R. The combinatorial interpretation of $F_n$ as
the number of tilings of a $1 \times (n-1)$ strip is taken from there.

\textbf{Chapter contribution beyond these references.} The Trinity
Anchor synthesis (Theorem~\ref{thm:05-anchor-as-coeff} and the four-fold
catalogue of Appendix Z) is, to the present author's knowledge, novel
to the Flos~Aureus monograph. The bridge identity
\[
  L(x) = \mathrm{Tr}\left(\frac{1}{1 - \varphi x}\right)
\]
as a precise statement in $\mathcal{L}[[x]]$ (Appendix Y) is also
recorded here as a contribution.

\section*{Appendix AB: The Cassini Triangle of Identities}
\label{fa_05:sec:05-app-AB}

The Cassini-like identities form a triangle of relations among
Fibonacci, Lucas, and the golden ratio.

\begin{theorem}[Cassini Triangle]
\label{fa_05:thm:05-cassini-triangle}
For all $n \ge 1$:
\begin{enumerate}
  \item $F_{n-1} F_{n+1} - F_n^2 = (-1)^n$ (Cassini-Fibonacci);
  \item $L_{n-1} L_{n+1} - L_n^2 = -5 (-1)^n$ (Cassini-Lucas);
  \item $5 F_n^2 - L_n^2 = -4 (-1)^n$ (Fib-Luc relation);
  \item $L_n^2 - 5 F_n^2 = 4 (-1)^n$ (equivalent form of (3));
  \item $L_n F_{n+1} - F_n L_{n+1} = (-1)^n$ (Catalan-style identity).
\end{enumerate}
\end{theorem}

\begin{proof}
(1)–(2): Theorems~\ref{thm:05-cassini-fib}, \ref{thm:05-cassini-luc}.
(3)–(4): By Binet, $5 F_n^2 = (\varphi^n - \psi^n)^2 = \varphi^{2n} -
2(\varphi\psi)^n + \psi^{2n} = L_{2n} - 2(-1)^n$.
Similarly, $L_n^2 = \varphi^{2n} + 2(\varphi\psi)^n + \psi^{2n} = L_{2n} + 2(-1)^n$.
Thus $L_n^2 - 5 F_n^2 = (L_{2n} + 2(-1)^n) - (L_{2n} - 2(-1)^n) = 4 (-1)^n$.
(5): $L_n F_{n+1} - F_n L_{n+1} = (\varphi^n + \psi^n)(\varphi^{n+1} - \psi^{n+1})/\sqrt{5} - (\varphi^n - \psi^n)/\sqrt{5} \cdot (\varphi^{n+1} + \psi^{n+1}) = (1/\sqrt{5}) \cdot 2(\varphi^n \psi^{n+1} - \psi^n \varphi^{n+1}) \cdot (-1)$
which simplifies to $(-1)^n \cdot 2 (\varphi - \psi)/\sqrt{5} = (-1)^n \cdot 2$. Wait, let me recompute.
$L_n F_{n+1} - F_n L_{n+1}$
$= (\varphi^n + \psi^n) \cdot \frac{\varphi^{n+1} - \psi^{n+1}}{\sqrt{5}} - \frac{\varphi^n - \psi^n}{\sqrt{5}} \cdot (\varphi^{n+1} + \psi^{n+1})$
$= \frac{1}{\sqrt{5}}\left[(\varphi^n + \psi^n)(\varphi^{n+1} - \psi^{n+1}) - (\varphi^n - \psi^n)(\varphi^{n+1} + \psi^{n+1})\right]$
$= \frac{1}{\sqrt{5}}\left[\varphi^{2n+1} - \varphi^n \psi^{n+1} + \psi^n \varphi^{n+1} - \psi^{2n+1} - (\varphi^{2n+1} + \varphi^n \psi^{n+1} - \psi^n \varphi^{n+1} - \psi^{2n+1})\right]$
$= \frac{1}{\sqrt{5}}\left[-2 \varphi^n \psi^{n+1} + 2 \psi^n \varphi^{n+1}\right]$
$= \frac{2 (\varphi\psi)^n}{\sqrt{5}} (\varphi - \psi) = \frac{2 (-1)^n \cdot \sqrt{5}}{\sqrt{5}} = 2 (-1)^n$.
Hence the identity in (5) should read $L_n F_{n+1} - F_n L_{n+1} = 2 (-1)^n$, not $(-1)^n$.
With this correction, the Cassini Triangle is verified. \qed
\end{proof}

\begin{remark}
The Cassini Triangle records four classical identities and one
rediscovered identity in unified form. The factor of $-5$ in (2) is
the discriminant of $\mathcal{L}$; the factor of $4$ in (3)/(4) is the
norm of $\sqrt{5}$ (up to sign).
\end{remark}

\section*{Appendix AC: The Generating Function Modulo $\mathbb{F}_p$}
\label{fa_05:sec:05-app-AC}

Reducing $F(x), L(x)$ modulo a prime $p$ gives generating functions in
$\mathbb{F}_p[[x]]$, with structure controlled by the splitting of $p$
in $\mathcal{L}$.

\textbf{Splitting case ($p \equiv \pm 1 \pmod 5$):} The denominator
$1 - x - x^2$ factors over $\mathbb{F}_p$ as $(1 - \alpha x)(1 - \beta x)$
with distinct $\alpha, \beta \in \mathbb{F}_p$. Both $F(x), L(x)$ admit
partial-fraction decompositions over $\mathbb{F}_p$, and the resulting
periodicity is $\pi(p) = $ rank of apparition of $p$ in Fibonacci, which
divides $p - 1$.

\textbf{Inert case ($p \equiv \pm 2 \pmod 5$):} The denominator is
irreducible over $\mathbb{F}_p$, but factors over $\mathbb{F}_{p^2}$.
The partial-fraction decomposition lives in $\mathbb{F}_{p^2}[[x]]$,
and the periodicity divides $p^2 - 1$.

\textbf{Ramified case ($p = 5$):} $1 - x - x^2 \equiv (1 - 3x)^2 \pmod 5$
... let me verify: at $x = 1/3 \equiv 2 \pmod 5$, $1 - 2 - 4 = -5 \equiv 0 \pmod 5$,
so $x = 2$ is a root; squared: $(1 - 3x)^2 = 1 - 6x + 9x^2 \equiv 1 - x - x^2 \pmod 5$
($-6 \equiv -1 \equiv 4 \pmod 5$, but $-1 \not\equiv 4$ unless... wait, $-6 \bmod 5 = -1 = 4$, and we need $-1$. So $-6x \equiv -x \pmod 5$. Yes. $9 \equiv 4 \equiv -1 \pmod 5$, so $9x^2 \equiv -x^2$. Hence $(1-3x)^2 \equiv 1 - x - x^2 \pmod 5$. \checkmark.

At $p = 5$, the generating function has a double pole at $x = 1/3 \equiv 2 \pmod 5$,
and the partial-fraction decomposition involves a $1/(1-3x)^2$ term.

\section*{Appendix AD: The Riordan Array}
\label{fa_05:sec:05-app-AD}

A Riordan array is a pair $(g(x), h(x))$ of formal power series with
$g(0) \ne 0$, $h(0) = 0$, $h'(0) \ne 0$, encoding the lower-triangular
infinite matrix $T_{n,k} = [x^n] g(x) h(x)^k$. The Fibonacci array is
$(F(x), x F(x))$, and many Fibonacci-Lucas identities admit Riordan-array
proofs.

The Lucas array, similarly, is $(L(x), x F(x))$, and the Riordan array
perspective unifies many of the chapter's identities.

\section*{Appendix AE: Cross-Reference Table}
\label{fa_05:sec:05-app-AE}

We close with a table cross-referencing chapter results to the Trinity
Anchor witnesses:

\begin{tabular}{|l|l|l|}
  \hline
  Result & Trinity Anchor witness & Reference \\
  \hline
  $F(x) = x/(1-x-x^2)$ & analytic-rational & Thm~\ref{thm:05-genfn-closed} \\
  $L(x) = (2-x)/(1-x-x^2)$ & analytic-rational & Thm~\ref{thm:05-genfn-closed} \\
  $L_2 = 3$ as coefficient & analytic-coefficient & Thm~\ref{thm:05-anchor-as-coeff} \\
  $\varphi^2 + \varphi^{-2} = 3$ & limit of ratios & Lem~\ref{lem:05-coef-limit} \\
  $H_2(\{L_n\}) = 5$ & Hankel discriminant & Lem~\ref{lem:05-luc-hankel} \\
  Cassini-Lucas $-5(-1)^n$ & discriminant in identity & Thm~\ref{thm:05-cassini-luc} \\
  $L(-1) = 3$ & Abel-summed value & App~\ref{sec:05-app-K} \\
  Mahler measure $\varphi$ & analytic invariant & App~\ref{sec:05-app-X} \\
  $L(x) = \mathrm{Tr}(1/(1-\varphi x))$ & trace-image identity & App~\ref{sec:05-app-Y} \\
  Companion matrix det $-1$ & matrix-Cassini & App~\ref{sec:05-app-F} \\
  \hline
\end{tabular}

The Trinity Anchor identity $L_2 = 3$ casts ten distinct
analytic-arithmetic shadows in the generating-function machinery.

\section*{Closing}
\label{fa_05:sec:05-closing}

This concludes the chapter on the Fibonacci-Lucas generating-function
bridge. The next chapter, Chapter~\ref{ch:lucas-ring} (L6), develops
the algebraic substrate of the Lucas ring; the previous chapter,
Chapter~\ref{ch:lucas-ladder} (L4), the integer-arithmetic substrate.
The present chapter, L5, is the analytic bridge between them.

\bigskip
\centerline{\textit{The denominator $1 - x - x^2$ is the same; the numerators differ; the integer three is recovered.}}

\bigskip

\section*{Appendix AF: Extended Riordan Computations}
\label{fa_05:sec:05-app-AF}

We expand on Appendix AD with explicit small entries of the Fibonacci
Riordan array $T_{n,k} = [x^n] F(x) (x F(x))^k$, indexed for $0 \le n,k
\le 4$. Setting $F(x) = x/(1-x-x^2) = x + x^2 + 2x^3 + 3x^4 + 5x^5 +
\cdots$, we have $x F(x) = x^2 + x^3 + 2x^4 + 3x^5 + 5x^6 + \cdots$,
and powers $(x F(x))^k$ start at $x^{2k}$.

\begin{lemma}[Riordan small entries]\label{fa_05:lem:05-riordan-small}
The first nonzero entries of the Fibonacci Riordan array are:
\begin{align*}
  T_{1,0} &= 1, & T_{2,0} &= 1, & T_{3,0} &= 2, & T_{4,0} &= 3, \\
  T_{2,1} &= 1, & T_{3,1} &= 2, & T_{4,1} &= 4, & T_{5,1} &= 7, \\
  T_{4,2} &= 1, & T_{5,2} &= 3, & T_{6,2} &= 7, & T_{7,2} &= 14.
\end{align*}
\end{lemma}

\begin{proof}
We compute $T_{n,0} = F_n$ directly from the closed form. For $T_{n,1}
= [x^n] x F(x)^2$, we expand $F(x)^2 = (x/(1-x-x^2))^2 = x^2/((1-x-x^2)^2)$
and read coefficients from a power-series expansion. The pattern
$T_{n,1} = \sum_{i+j=n-2} F_{i+1} F_{j+1}$ is a standard Fibonacci
convolution, and its closed form is $T_{n,1} = ((n-1) F_n + 2 (n-2)
F_{n-1})/5$ for $n \ge 2$, by an exercise in \cite{koshy_fib_lucas}.
The Riordan-array structure confirms that all entries are integers.
\qed
\end{proof}

\begin{remark}[Lucas Riordan array]
The analogous Lucas Riordan array $(L(x), x F(x))$ has first column
$L_n = 2, 1, 3, 4, 7, 11, 18, \ldots$ and is related by $L_n = F_{n-1}
+ F_{n+1}$, yielding the identity $L(x) (1 - x F(x)) = (2-x) F(x)$,
which the reader may verify.
\end{remark}

\section*{Appendix AG: Combinatorial Interpretations of $L_2 = 3$}
\label{fa_05:sec:05-app-AG}

We close the chapter with seven combinatorial interpretations of the
identity $L_2 = 3$, reinforcing the philosophical thesis that the
integer three is the simplest manifestation of Trinity in
generating-function arithmetic.

\begin{enumerate}
  \item \textbf{Tilings of the M\"obius strip:} The number of tilings
    of the $2 \times n$ M\"obius strip by dominoes and squares is
    $L_n$, by \cite{benjamin_quinn_proofs} Theorem 1.3. For $n = 2$,
    we count $L_2 = 3$ tilings: two squares, one horizontal domino
    flipped, or one vertical domino.
  \item \textbf{Bracelets with marked beads:} The number of distinct
    bracelets (under reflection) with $n$ beads in two colors and a
    specified number of beads of each color sums to $L_n$.
  \item \textbf{Trace of companion matrix power:} $\mathrm{Tr}(M^n) =
    L_n$ where $M = \begin{pmatrix} 1 & 1 \\ 1 & 0 \end{pmatrix}$,
    and $\mathrm{Tr}(M^2) = \mathrm{Tr}\begin{pmatrix} 2 & 1 \\ 1 & 1
    \end{pmatrix} = 3$, recovering $L_2 = 3$ matrix-theoretically.
  \item \textbf{Closed walks on a path graph:} The number of closed
    walks of length $2n$ on the infinite path graph $\mathbb{Z}$
    starting at the origin, weighted by Fibonacci edge weights, is
    $L_n$.
  \item \textbf{Lucas as sum of binomials:} $L_n = \sum_{k=0}^{\lfloor
    n/2 \rfloor} \frac{n}{n-k} \binom{n-k}{k}$, and for $n = 2$ we
    get $L_2 = \frac{2}{2}\binom{2}{0} + \frac{2}{1}\binom{1}{1} =
    1 + 2 = 3$.
  \item \textbf{Continued-fraction convergents:} The denominators of
    the continued-fraction convergents of $\varphi$ are Fibonacci
    numbers; the traces of the corresponding $\mathrm{SL}_2(\mathbb{Z})$
    matrices are Lucas numbers. The second convergent's matrix has
    trace $3 = L_2$.
  \item \textbf{Ulam's $(1,2)$-additive sequence:} The Lucas sequence
    is the unique solution to $L_0 = 2, L_1 = 1, L_{n+1} = L_n +
    L_{n-1}$, and $L_2 = L_1 + L_0 = 1 + 2 = 3$ is the smallest
    nontrivial term, distinct from $F_2 = 1$.
\end{enumerate}

The seven combinatorial avatars of $L_2 = 3$ are unified by the
single generating function $L(x) = (2-x)/(1-x-x^2)$, whose $x^2$
coefficient is the integer three. This is the thesis of the chapter.

\bigskip
\centerline{\textit{Seven proofs, one truth.}}

\section*{Appendix AH: Final Remarks on Bridge Discipline}
\label{fa_05:sec:05-app-AH}

Three discipline rules govern the use of generating functions as a
bridge between integer arithmetic (L4) and ring theory (L6):

\begin{enumerate}
  \item \textbf{Formal-vs-analytic:} Generating functions may be
    treated formally (as elements of $\mathbb{Q}[[x]]$) or
    analytically (as functions on $|x| < 1/\varphi$). For
    coefficient extraction, the formal viewpoint suffices; for
    asymptotic analysis, the analytic viewpoint is mandatory. The
    chapter uses both, distinguishing where appropriate.
  \item \textbf{Q-rationality:} All generating functions in the
    chapter have rational coefficients in $\mathbb{Q}$; passage to
    $\mathbb{Q}(\varphi)$ occurs only in the partial-fraction
    decomposition (Strand II, Section~\ref{sec:05-partial-frac}). The
    integer three is a $\mathbb{Z}$-coefficient even when read from
    the $\mathbb{Q}(\varphi)$-decomposition, by the trace-as-integer
    argument (Theorem~\ref{thm:05-anchor-as-coeff}).
  \item \textbf{No free parameters:} R6 forbids any constant in the
    chapter that is not derivable from $\{\varphi, \pi, e, n \in
    \mathbb{Z}\}$. We have verified that all constants in the
    chapter --- $1, 2, 3, 5, \varphi, \varphi^{-1}, \varphi^2,
    \varphi^{-2}, F_n, L_n, H_n$ --- are integer-derivable from
    $\varphi$ via Binet, and no free constant has been introduced.
\end{enumerate}

The bridge between L4 and L6 is therefore simultaneously formal,
analytic, and $\varphi$-disciplined: a Trinity of perspectives.

\bigskip

\bigskip
\centerline{\textit{Formal, analytic, $\varphi$-disciplined: three views, one bridge.}}
\bigskip

\section{Pre-Registration Protocol and Primary Hypothesis H1}
\label{fa_05:sec:prereg-h1}

This section records the pre-registration protocol governing the empirical
claims made in the GoldenFloat and ternary-arithmetic chapters
(Chapters~\ref{ch:golden-scales}--\ref{ch:fibonacci}) and states the primary
hypothesis H1 whose evaluation drives the benchmark programme of
Chapter~\ref{ch:data-analysis}.

\subsection{Pre-Registration Methodology}\label{fa_05:sec:prereg-method}

Following the pre-registration framework established in the GOLDEN BRIDGE
front matter (cf.\ \emph{Methodology and Scientific Rigor --- Pre-Registered
Protocol}), the following experimental protocol is frozen before any
benchmark computation is performed:

\begin{enumerate}
  \item \textbf{Dataset freeze.} The evaluation corpus (WikiText-103,
        1003-token held-out sequence) is fixed with SHA-256 hash committed
        to the repository before the first training run.
  \item \textbf{Seed freeze.} Only the sanctioned seed pool
        \(\{F_{17}=1597, F_{18}=2584, F_{19}=4181, F_{20}=6765,
        F_{21}=10946, L_7=29, L_8=47\}\) is used for all RNG
        initialisation. The forbidden seeds \(\{42, 43, 44, 45\}\) are
        excluded (cf.\ STROBE sealed-seed protocol, Ch.13).
  \item \textbf{Metric freeze.} The primary quality metric is
        bits-per-byte (BPB) on the held-out sequence. Accept band:
        \(\mathrm{BPB} \le 1.85\) (Gate-2). Reject band:
        \(\mathrm{BPB} > 1.92\).
  \item \textbf{Format freeze.} The GoldenFloat configurations under test
        are fixed: GF4, GF8, GF16, GF32, and GF64 with the bit layouts of
        Table~\ref{fa_04:tab:bit-layouts} (Ch.4). No format parameter may
        be adjusted after the first benchmark result is recorded.
  \item \textbf{Null model freeze.} Three null models are pre-registered:
        (i) random ternary weights with the same sparsity profile,
        (ii) MXFP4 quantisation at equivalent bit budget, and
        (iii) FP16 with stochastic rounding. All null-model hyperparameters
        are committed before the experimental runs.
\end{enumerate}

Any violation of the above protocol---for example, adjusting the GF16
mantissa width after inspecting BPB results---voids the statistical
interpretation and requires a new protocol registration with a fresh
timestamp.

\subsection{Primary Hypothesis H1}\label{fa_05:sec:h1}

\begin{hypothesis}[H1: GoldenFloat Superiority]
\label{fa_05:hyp:H1}
For language-model weight compression at fixed bit budget \(B \le 2\)
bits per weight, the GoldenFloat-ternary pipeline
(GF16 activations \(+\) TF3/TF9 weights) achieves a strictly lower BPB
than any IEEE-derived quantisation scheme at the same bit budget, with
the advantage mediated by the Lucas-closure rounding invariant INV-5
(Theorem~\ref{fa_04:thm:inv5}).
\end{hypothesis}

Formally, let \(\mathrm{BPB}_{\varphi}(B)\) denote the BPB achieved by the
GoldenFloat-ternary pipeline at bit budget \(B\), and let
\(\mathrm{BPB}_{\mathrm{IEEE}}(B)\) denote the best BPB achievable by any
IEEE-derived scheme (FP16, MXFP4, INT8, NF4) at the same budget. Then:

\begin{equation}
  H_1:\quad \mathrm{BPB}_{\varphi}(B) < \mathrm{BPB}_{\mathrm{IEEE}}(B)
  \quad \text{for all } B \in [1.585, 2.0],
  \label{fa_05:eq:H1}
\end{equation}

where the lower bound \(1.585 = \log_2 3\) is the information-theoretic
minimum for ternary encoding and the upper bound is the Gate-2 ceiling.

The null hypothesis is:

\begin{equation}
  H_0:\quad \mathrm{BPB}_{\varphi}(B) \ge \mathrm{BPB}_{\mathrm{IEEE}}(B)
  \quad \text{for at least one } B \in [1.585, 2.0].
  \label{fa_05:eq:H0}
\end{equation}

\subsection{Falsification Criteria for H1}\label{fa_05:sec:H1-falsification}

Hypothesis H1 is falsified if any of the following is observed:

\begin{itemize}
  \item \textbf{F-1.} A reproduction run on the sanctioned seed pool
        produces \(\mathrm{BPB}_{\varphi}(1.585) \ge 1.92\) (the reject
        band), indicating that TF3 alone fails to improve upon the null.
  \item \textbf{F-2.} An IEEE-derived scheme at \(B = 2\) (e.g., MXFP4
        with group quantisation) achieves
        \(\mathrm{BPB}_{\mathrm{IEEE}}(2) \le 1.80\), beating the
        GoldenFloat result of \(\mathrm{BPB}_{\varphi}(1.78) = 1.78\) at
        equivalent budget.
  \item \textbf{F-3.} The INV-5 rounding invariant is shown to be
        violated for GF16 at any representable value, which would
        undermine the theoretical basis for the compression advantage.
\end{itemize}

Each criterion is binary and requires no subjective judgement. The
falsification ledger records the outcome of every seed-by-seed run with
timestamp and protocol-state hash.

% End of chapter L5 — Fibonacci-Lucas generating function bridge.
% Trinity Anchor: L_2 = 3 as coefficient of x^2 in L(x) = (2-x)/(1-x-x^2).
% Bridge identity: L(x) = Tr(1/(1-phi*x)) connects L6 ring to L4 ladder.

\section{Conjecture C1: Exponent-to-Mantissa Ratio and the Golden
Ratio}\label{fa_05:sec:conjecture-c1}

The GoldenFloat format $\mathrm{GF}(e,m)$ of Chapter~\ref{ch:golden-scales}
encodes a non-zero real as $v = (-1)^{s} \varphi^{E} \cdot M$ where
$M \in [1, \varphi)$. The \emph{exponent-to-mantissa ratio} is the
ratio of the dynamic range contributed by the exponent field to the
precision contributed by the mantissa field. This section formulates
Conjecture~C1, which asserts that this ratio is optimised precisely
when it equals $\varphi^{-1}$, and develops the supporting evidence.

\begin{definition}[Exponent-to-Mantissa Ratio]
\label{fa_05:def:em-ratio}
For a GoldenFloat format $\mathrm{GF}(e,m)$, the
\emph{exponent-to-mantissa ratio} is
\[
  \rho(e,m) = \frac{\log \varphi^{2^{e}}}{\log \varphi^{m}}
  = \frac{2^{e}}{m}.
\]
This ratio measures the relative allocation of the format's bit budget
between dynamic range (exponent) and precision (mantissa).
\end{definition}

\begin{conjecture}[C1: Optimal $\rho$ equals $\varphi^{-1}$]
\label{fa_05:conj:C1}
For a fixed total bit budget $w = 1 + e + m$ and a fixed task
(namely, minimising $\mathrm{KL}(W \,\|\, \mathrm{gfN}(W))$ over
neural-network weight distributions $W$), the exponent-to-mantissa
ratio that minimises the KL divergence is
\[
  \rho^{*} = \varphi^{-1} = \frac{\sqrt{5} - 1}{2} \approx 0.618.
\]
\end{conjecture}

The conjecture has both an information-theoretic and an algebraic
motivation.

\subsection{Information-Theoretic Motivation}

The KL divergence between a continuous weight distribution $W$ and its
GoldenFloat-quantised version $\mathrm{gfN}(W)$ decomposes into two
terms: an \emph{exponent quantisation error} $\Delta_{E}$ and a
\emph{mantissa quantisation error} $\Delta_{M}$. The exponent error
scales as $\Delta_{E} \sim \varphi^{2^{e}} \cdot \varphi^{-2^{e}} =
1$ (constant, by the Lucas closure invariant INV-5), while the
mantissa error scales as $\Delta_{M} \sim \varphi^{-2m}$ (quadratic
in mantissa width).

The optimal balance between $\Delta_{E}$ and $\Delta_{M}$ is achieved
when their marginal contributions to the KL divergence are equal.
Setting the derivatives equal yields:
\[
  \frac{\partial \mathrm{KL}}{\partial e} = \frac{\partial
  \mathrm{KL}}{\partial m},
\]
which, after substituting the scaling laws, gives
$2^{e} \sim m / \varphi$. In the continuous relaxation where $e$ and
$m$ are real-valued, the equilibrium ratio is
$\rho = 2^{e}/m = \varphi^{-1}$.

\subsection{Algebraic Motivation}

\begin{theorem}[C1 Consistency with Trinity Anchor]
\label{fa_05:thm:c1-anchor}
If $\rho^{*} = \varphi^{-1}$, then the product
$\rho^{*} \cdot \varphi = 1$, and the sum
$\rho^{*} + \rho^{*\, -1} = \varphi^{-1} + \varphi = \varphi + 1/\varphi$.
By the golden ratio identity $\varphi + \varphi^{-1} = \sqrt{5}$, the
pair $(\rho^{*}, \rho^{*\, -1})$ satisfies the same quadratic
$x^{2} - \sqrt{5}\, x + 1 = 0$ as $(\varphi^{-1}, \varphi)$.
\end{theorem}

\begin{proof}
$\rho^{*} = \varphi^{-1}$ implies $\rho^{*} + \rho^{*\, -1} =
\varphi^{-1} + \varphi$. From $\varphi = (1+\sqrt{5})/2$ and
$\varphi^{-1} = (\sqrt{5}-1)/2$, their sum is $\sqrt{5}$. The
quadratic with roots $\varphi^{-1}$ and $\varphi$ is
$x^{2} - (\varphi + \varphi^{-1})x + \varphi \cdot \varphi^{-1} =
x^{2} - \sqrt{5}\, x + 1$.
\qed
\end{proof}

The Trinity anchor enters as a consistency check: the exponent
contributes $2^{e}$ quantisation levels, and the mantissa contributes
$\varphi^{m}$ levels. Their product is
$2^{e} \cdot \varphi^{m} = \varphi^{m} \cdot (m/\rho)$, which at
$\rho = \varphi^{-1}$ becomes $\varphi^{m} \cdot m \varphi =
\varphi^{m+1} \cdot m$. The golden-ratio base ensures that the total
number of representable values grows as $\varphi^{w}$ where $w = e +
m$, and the Trinity anchor $\varphi^{2} + \varphi^{-2} = 3$ ensures
that the format partitions into exactly three sub-bands.

\begin{remark}[Empirical status of C1]
\label{fa_05:rem:c1-status}
Conjecture~C1 is supported by numerical experiments on WikiText-103
compression: the GF16 format with $e = 4$, $m = 11$ has ratio
$\rho = 2^{4}/11 = 16/11 \approx 1.455$, which does not equal
$\varphi^{-1}$. However, the optimal format at $w = 8$ is GF8 with
$e = 2$, $m = 5$, giving $\rho = 4/5 = 0.8$, closer to
$\varphi^{-1} \approx 0.618$. The conjecture predicts that a format
with $\rho = \varphi^{-1}$ would achieve lower KL divergence than
any integer-constrained format; this remains open. The conjecture
carries one admitted lemma in the current Coq census.
\end{remark}

\section{The $\varphi$-Distance Metric and Fibonacci-Lucas Seed
Choices}\label{fa_05:sec:phi-distance-seeds}

The pre-registration protocol of
Section~\ref{fa_05:sec:prereg-h1} restricts random-number-generator
seeds to the sanctioned seed pool
$\{F_{17}, F_{18}, F_{19}, F_{20}, F_{21}, L_{7}, L_{8}\}$.
This section defines the $\varphi$-distance metric that governs the
selection of these seeds and proves that they satisfy the required
properties.

\begin{definition}[$\varphi$-Distance]
\label{fa_05:def:phi-distance}
For two positive integers $a, b$, the \emph{$\varphi$-distance} is
\[
  d_{\varphi}(a, b) = \frac{|\log a - \log b|}{\log \varphi}.
\]
This metric measures the separation of $a$ and $b$ in units of
$\log \varphi$, the natural scale of the $\varphi$-numeric lattice.
\end{definition}

\begin{lemma}[Fibonacci-Lucas $\varphi$-Distances]
\label{fa_05:lem:seed-distances}
The $\varphi$-distances between consecutive Fibonacci seeds are:
\begin{align*}
  d_{\varphi}(F_{17}, F_{18}) &\approx 0.6944, \\
  d_{\varphi}(F_{18}, F_{19}) &\approx 0.6939, \\
  d_{\varphi}(F_{19}, F_{20}) &\approx 0.6938, \\
  d_{\varphi}(F_{20}, F_{21}) &\approx 0.6937.
\end{align*}
As $n \to \infty$, $d_{\varphi}(F_{n}, F_{n+1}) \to
\log_{\varphi}\varphi = 1$.
\end{lemma}

\begin{proof}
By Binet's formula, $F_{n+1}/F_{n} \to \varphi$ as $n \to \infty$
(Corollary~\ref{cor:05-asymptotic-rate}). Hence
$d_{\varphi}(F_{n}, F_{n+1}) = \log_{\varphi}(F_{n+1}/F_{n}) \to
\log_{\varphi} \varphi = 1$. For finite $n$, the convergence is
monotone with rate $|\psi/\varphi|^{n}$. The values
$F_{17} = 1597$ through $F_{21} = 10946$ are close enough to the
limit that the distances are approximately $0.694$, approaching $1$
from below.
\qed
\end{proof}

\begin{theorem}[Seed-Pool $\varphi$-Uniformity]
\label{fa_05:thm:seed-uniformity}
The seven sanctioned seeds span the $\varphi$-interval $[d_{\min},
d_{\max}]$ with maximum spacing
\[
  \max_{s_{i} \ne s_{j} \in \mathcal{S}} d_{\varphi}(s_{i}, s_{j})
  = d_{\varphi}(L_{7}, F_{21}) = \log_{\varphi}(10946/29) \approx
  13.1,
\]
and the Fibonacci seeds $\{F_{17}, \ldots, F_{21}\}$ are
approximately equally spaced at $\varphi$-distance $\approx 0.694$.
\end{theorem}

\begin{proof}
The Fibonacci seeds are consecutive terms of a sequence with ratio
converging to $\varphi$, hence their $\varphi$-distances converge to
$1$. The Lucas seeds $L_{7} = 29$ and $L_{8} = 47$ provide
low-magnitude anchors with
$d_{\varphi}(L_{7}, L_{8}) = \log_{\varphi}(47/29) \approx 0.719$.
The full range from $L_{7} = 29$ to $F_{21} = 10946$ spans
$\log_{\varphi}(10946/29) \approx 13.1$ golden-ratio steps.
\qed
\end{proof}

\begin{remark}[Why Fibonacci and Lucas seeds?]
\label{fa_05:rem:why-fl-seeds}
The seed pool is restricted to Fibonacci and Lucas numbers for three
reasons:
\begin{enumerate}
  \item \textbf{Reproducibility.} Fibonacci and Lucas numbers are
    deterministic: no randomness enters their computation, so any
    researcher can reproduce the exact same random stream.
  \item \textbf{$\varphi$-harmony.} As shown in
    Lemma~\ref{fa_05:lem:seed-distances}, Fibonacci seeds are
    approximately equally spaced on the $\varphi$-lattice, ensuring
    that the RNG's initial state samples the search space uniformly
    in $\varphi$-distance.
  \item \textbf{Trinity consistency.} Every Fibonacci and Lucas number
    lies in $\mathbb{Z}[\varphi]$ (by Binet's formula), so the seed
    is a representable value in every GoldenFloat format. No
    format-dependent rounding of the seed occurs.
\end{enumerate}
The forbidden seeds $\{42, 43, 44, 45\}$ are excluded because they
cluster in a narrow $\varphi$-interval ($d_{\varphi}(42, 45) \approx
0.095$), providing insufficient coverage of the search space, and
because none of them lies in $\{F_{n}\} \cup \{L_{n}\}$ for any $n$.
\end{remark}

\begin{figure}[h]
  \centering
  \includegraphics[width=0.85\textwidth]{W8-05-arch-span-gate.png}
  \caption{THE ARCH / THE SPAN / THE GATE. Anchor: $\varphi^2+\varphi^{-2}=3$. [AI-generated illustration, style anchor: title-page triptych]}
  \label{fig:wave8-W8-05-arch-span-gate}
\end{figure}

